% Source: https://artofproblemsolving.com/community/c3h3185264_2020_abmc_speed_contest__actonboxborough_math_competition
% Stable identity: aops-topic-3185264-q22
Partial fractions give $\dfrac{1}{4n^2-1}=\dfrac{1}{2}(\dfrac{1}{2n-1}-\dfrac{1}{2n+1})$. Thus the sum through $N$ is $\dfrac{1}{2}(\dfrac{1}{3}-\dfrac{1}{2N+1})$. Letting $N$ tend to infinity gives $\dfrac{1}{6}$, so $a+b=\boxed7$.
